\section{Relation to symmetry, reciprocity and multi-speed identification}
Reflection symmetry tests forbidden even/odd components and gives a closest
pair under a declared metric. It does not supply the independent torque shape
used here, and the present objective does not re-identify a winding resistance.
Quantities normalized under a sole resistance-mismatch hypothesis remain
\emph{resistance-equivalent} rather than physical resistance estimates.

Reciprocity asks whether a field can be the gradient of one co-energy surface.
The selected numerical correction satisfies that model structurally, but the
raw terminal reconstruction may contain nonconservative measurement effects.
Adding a curl-free correction cannot cancel their curl. Low torque residual,
low curl and smoothness are separate admissibility statements, not mutually
validating proofs of true flux.

Multi-speed data can test whether a proposed correction transfers across
conditions. A resistance-induced reconstruction error ideally scales with
$1/\omega_e$, whereas an ideal static magnetic map does not. Current-proportional
voltage error has the same column as resistance mismatch at every speed, and
terminal current need not equal the magnetic storage current in a loss-aware
model. The companion identifiability paper owns the scaled rank, nuisance and
support analysis. This paper uses its conclusion operationally: a correction
that helps one speed but harms another is evidence against a common magnetic
interpretation under the current nuisance model.
